% SEGMENT OLP-0417-S01
% Part: normal-modal-logic
% Chapter: syntax-and-semantics
% Section: schemas

\documentclass[../../../include/open-logic-section]{subfiles}

\begin{document}

\olfileid{nml}{syn}{sch}

\olsection{திட்டவடிவங்களும் செல்லுபடித்தன்மையும்}

\begin{defn}
  ஏதோ ஒரு வாய்ப்புநிலை !!{formula}~$!C$ இன் பதிலீட்டு நிகழ்வுகள்
  அனைத்தையும், அவற்றை மட்டுமே கொண்ட !!{formula}s இன் கணமே ஒரு
  \emph{திட்டவடிவம்}; அதாவது,
  \[
  \Setabs{!B}{\lexists[!D_1], \dots, \lexists[!D_n] \left(!B =
  \SSubst{!C}{\subst{!D_1}{p_1}, \dots, \subst{!D_n}{p_n}} \right) }.
  \]
  !!{formula}~$!C$ ஆனது அந்தத் திட்டவடிவத்தின் \emph{பண்புறுத்தும்}
  !!{formula} எனப்படுகிறது; !!{propositional variable}s இன் பெயர்மாற்றம்
  வரையில் அது தனித்தது. !!^a{formula}~$!A$ அந்தக் கணத்தின் உறுப்பாக
  இருந்தால், அது ஒரு திட்டவடிவத்தின் \emph{நிகழ்வு}.
\end{defn}

$!C$ இன் அணுக் கூறுகளுக்குப் பதிலாக `$!A$', `$!B$',~\dots, ஆகியவற்றை
இடுவதால் கிடைக்கும் மேல்மொழிக் கோவையால் ஒரு திட்டவடிவத்தைக் குறிப்பது
வசதியானது. ஆகவே, எடுத்துக்காட்டாக, பின்வருவன திட்டவடிவங்களைக் குறிக்கின்றன:
`$!A$', `$!A \lif \Box!A$', `$!A \lif (!B \lif !A)$'. அவை முறையே
$p$, $p \lif \Box p$, $p \lif (q \lif p)$ என்ற பண்புறுத்தும்
!!{formula}s க்கு ஒத்தவை. `$!A$' என்ற திட்டவடிவம் \emph{எல்லா}
!!{formula}s இன் கணத்தையும் குறிக்கிறது.

\begin{defn}
  ஒரு திட்டவடிவத்தின் எல்லா நிகழ்வுகளும் ஒரு மாதிரியில் மெய்யாக இருந்தால்,
  எனில் மட்டுமே, அந்தத் திட்டவடிவம் அம்மாதிரியில் \emph{மெய்}; மேலும் ஒரு
  திட்டவடிவம் ஒவ்வொரு மாதிரியிலும் மெய்யாக இருந்தால், எனில் மட்டுமே, அது
  \emph{செல்லுபடியாகும்}.
\end{defn}

\begin{prop}\ollabel{prop:Kvalid}
  பின்வரும் திட்டவடிவம் \Ax{K} செல்லுபடியாகும்:
  \begin{equation*}
  \Box(!A \lif !B) \lif (\Box !A \lif \Box !B). \tag{\Ax{K}}
  \end{equation*}
\end{prop}

\begin{proof}
  திட்டவடிவத்தின் எல்லா நிகழ்வுகளும் ஒவ்வொரு மாதிரியின் ஒவ்வோர் உலகிலும்
  மெய்யாகின்றன என்று நாம் காட்ட வேண்டும். ஆகவே $\mModel{M} =
  \tuple{W,R,V}$ மற்றும் $w \in W$ ஆகியவை ஏதேனுமாக இருக்கட்டும். ஒரு
  நிபந்தனைக் கூற்று ஓர் உலகில் மெய்யாகும் என்று காட்ட, அதன் முன்கூறு
  மெய்யென்று கொண்டு அதன் பின்கூறும் மெய்யென்று காட்டுகிறோம். இந்நேர்வில்,
  $\mSat{M}{\Box(!A \lif !B)}[w]$ மற்றும் $\mSat{M}{\Box !A}[w]$ என்க.
  $\mSat{M}{\Box !B}[w]$ என்று நாம் காட்ட வேண்டும். எனவே $Rww'$ ஆகும்
  ஏதேனும் ஒரு $w'$ ஐக் கொள்க. முதல் கருதுகோளால்
  $\mSat{M}{!A \lif !B}[w']$; இரண்டாம் கருதுகோளால்
  $\mSat{M}{!A}[w']$. இதிலிருந்து $\mSat{M}{!B}[w']$ பின்வருகிறது. $w'$
  ஏதேனுமாக இருந்ததால், $\mSat{M}{\Box !B}[w]$.
\end{proof}

\begin{prop}\ollabel{prop:Dual-valid}
  பின்வரும் திட்டவடிவம் \Dual{} செல்லுபடியாகும்:
  \begin{equation*}
  \Diamond !A \liff \lnot\Box\lnot !A. \tag{\Dual}
  \end{equation*}
\end{prop}

\begin{proof}
  பயிற்சி.
\end{proof}

\begin{prob}
  \olref[nml][syn][sch]{prop:Dual-valid} ஐ நிறுவுக.
\end{prob}

\begin{prop}\ollabel{prop:soundMP}
  $!A$ மற்றும் $!A \lif !B$ ஆகியவை ஒரு மாதிரியின் ஓர் உலகில் மெய்யாக
  இருந்தால், $!B$ உம் அவ்வுலகில் மெய்யாகும். ஆகவே செல்லுபடியான
  !!{formula}s மோடஸ் போனென்ஸின் கீழ் மூடப்பட்டவை.
\end{prop}

\begin{prop}\ollabel{prop:valid-instances}
  !!^a{formula}~$!A$ செல்லுபடியாகும் என்பது, அப்போதும் அப்போது மட்டுமே,
  அதன் எல்லாப் பதிலீட்டு நிகழ்வுகளும் செல்லுபடியாகும். வேறு சொற்களில்,
  ஒரு திட்டவடிவம் செல்லுபடியாகும் என்பது, அப்போதும் அப்போது மட்டுமே, அதன்
  பண்புறுத்தும் !!{formula} செல்லுபடியாகும்.
\end{prop}

\begin{proof}
  $!A$ தனக்கே ஒரு பதிலீட்டு நிகழ்வு என்பதால், ``எனில்'' திசை தெளிவு.

  ``எனில் மட்டுமே'' திசையை நிறுவ, பின்வருவதைக் காட்டுகிறோம்:
  $\mModel{M} = \tuple{W, R, V}$ ஒரு வாய்ப்புநிலை மாதிரியாகவும்,
  $!B \ident
  \SSubst{!A}{\subst{!D_1}{p_1},\dots,\subst{!D_n}{p_n}}$ ஆனது~$!A$ இன்
  ஒரு பதிலீட்டு நிகழ்வாகவும் இருக்கட்டும். $V'(p_i) =
  \Setabs{w}{\mSat{M}{!D_i}[w]}$ என்று வைத்து $\mModel{M'} = \tuple{W,
  R, V'}$ ஐ வரையறுக்க. அப்போது எந்த~$w \in W$ க்கும்
  $\mSat{M}{!B}[w]$ எனில், எனில் மட்டுமே, $\mSat{M'}{!A}[w]$.
  (நிறுவலைப் பயிற்சியாக விடுகிறோம்.) இப்போது $!A$ செல்லுபடியாகவும்,
  $!A$ இன் ஏதோ ஒரு பதிலீட்டு நிகழ்வு $!B$ செல்லுபடியாகாமலும் இருப்பதாகக்
  கொள்க. அப்போது ஏதோ ஒரு
  $\mModel{M} = \tuple{W, R, V}$ மற்றும் ஏதோ ஒரு $w \in W$ க்கு
  $\mSat/{M}{!B}[w]$. ஆனால் மேற்கூற்றின்படி $\mSat/{M'}{!A}[w]$; எனவே
  $!A$ செல்லுபடியன்று. இது ஒரு முரண்பாடு.
\end{proof}

\begin{prob}
  \olref[nml][syn][sch]{prop:valid-instances} இன் நிறுவலிலுள்ள ``எனில்
  மட்டுமே'' பகுதியில் கூறப்பட்ட மேற்கூற்றை நிறுவுக. (குறிப்பு:~$!A$ மீது
  தொகுத்தறிதலைப் பயன்படுத்துக.)
\end{prob}

எனினும், ஒரு திட்டவடிவம் ஒரு மாதிரியில் மெய்யாகும் என்பது, அப்போதும்
அப்போது மட்டுமே, அதன் பண்புறுத்தும் வாய்பாடு அம்மாதிரியில் மெய்யாகும்
என்பது சரியன்று. ``எனில் மட்டுமே'' திசை மெய்யே: $!A$ இன் ஒவ்வொரு
நிகழ்வும்~$\mModel{M}$ இல் மெய்யாக இருந்தால், $!A$ தானும்~$\mModel{M}$
இல் மெய்யாகும். ஆனால் $!A$ ஆனது~$\mModel{M}$ இல் மெய்யாகவும், $!A$ இன்
ஏதோ ஒரு நிகழ்வு~$\mModel{M}$ இன் ஏதோ ஓர் உலகில் பொய்யாகவும் இருக்கலாம்.
மிக எளிய ஓர் எதிரெடுத்துக்காட்டாக, $w$ என்ற ஒரே ஓர் உலகையும்
$V(p) = \{w\}$ ஐயும் கொண்ட மாதிரியில்~$p$ ஐக் கருதுக; அதனால் $w$ இல்
$p$ மெய். ஆனால் $\lfalse$ ஆனது~$p$ இன் ஒரு நிகழ்வு; அது~$w$ இல்
மெய்யன்று.

\begin{prob}
  பின்வரும் !!{formula}s எதுவும் செல்லுபடியாகாது என்று காட்டுக:
  \begin{enumerate}
    \item[\Ax{D}:] \quad $\Box p \lif \Diamond p$;
    \item[\Ax{T}:] \quad $\Box p \lif p$;
    \item[\Ax{B}:] \quad $p \lif \Box\Diamond p$;
    \item[\Ax{4}:] \quad $\Box p \lif \Box \Box p$;
    \item[\Ax{5}:] \quad $\Diamond p \lif \Box \Diamond
      p$.
  \end{enumerate}
\end{prob}

\begin{table}[t]
    \centering
    \begin{tabular}{| l || l |}
      \hline
      {\emph{செல்லுபடியான திட்டவடிவங்கள்}} & {\emph{செல்லுபடியாகாத திட்டவடிவங்கள்}} \\
      \hline\hline
      $\Box(!A \lif !B) \lif (\Diamond !A \lif \Diamond !B)$
      & $\Box (!A \lor !B) \lif (\Box !A \lor \Box !B)$ \\
      $\Diamond (!A \lif !B) \lif (\Box !A \lif \Diamond
      !B)$
      & $(\Diamond !A \land \Diamond !B) \lif \Diamond (!A
      \land !B)$\\
      $\Box (!A \land !B) \liff (\Box !A \land \Box !B)$
      & $!A \lif \Box !A$ \\
      $\Box !A \lif \Box (!B \lif !A)$
      & $\Box \Diamond !A \lif !B$ \\
      $\lnot \Diamond !A \lif \Box (!A \lif !B)$
      & $\Box \Box !A \lif \Box !A$ \\
      $\Diamond (!A \lor !B) \liff (\Diamond !A \lor
      \Diamond !B)$
      & $\Box \Diamond !A \lif \Diamond \Box !A$. \\
      \hline
    \end{tabular}
    % SOURCE CORRECTION TA-NML-010: remove the source's editorial uncertainty; both columns are explicitly classified and exercised below.
    \caption{செல்லுபடியான மற்றும் செல்லுபடியாகாத திட்டவடிவங்கள்.}
    \ollabel{tab:valid-invalidSchemas}
\end{table}

\begin{prob}%\ollabel{ex:in/validSchemas}
  \olref[nml][syn][sch]{tab:valid-invalidSchemas} இன் முதல் நெடுக்கையில்
  உள்ள திட்டவடிவங்கள் செல்லுபடியாகின்றன என்றும், இரண்டாம் நெடுக்கையில்
  உள்ளவை செல்லுபடியாகவில்லை என்றும் நிறுவுக.
\end{prob}

\begin{prob}
  பின்வரும் திட்டவடிவங்கள் செல்லுபடியாகின்றனவா, செல்லுபடியாகவில்லையா என்று
  தீர்மானிக்க:
  \begin{enumerate}
  \item $(\Diamond !A \lif \Box !B) \lif (\Box !A \lif \Box
    !B)$;
  \item $\Diamond(!A \lif !B) \lor \Box(!B \lif !A)$.
  \end{enumerate}
\end{prob}

\begin{prob}
  பின்வரும் ஒவ்வொரு திட்டவடிவத்திற்கும், அந்த !!{formula} வின் ஒவ்வொரு
  நிகழ்வும் $\mModel{M}$ இல் மெய்யாகுமாறு ஒரு மாதிரி $\mModel{M}$ ஐக்
  காண்க:
  \begin{enumerate}
  \item $p \lif \Diamond\Diamond p$;
  \item $\Diamond p \lif \Box p$.
  \end{enumerate}
\end{prob}

\end{document}
